\documentclass{beamer}
\usepackage{beamerthemeshadow}
\usepackage[utf8]{inputenc}  
\usepackage{longtable}
\usepackage{minted}
\usepackage[ngermanb]{babel}
\beamersetuncovermixins{\opaqueness<1>{25}}{\opaqueness<2->{15}}
\setbeamertemplate{navigation symbols}{}%remove navigation symbols
\begin{document}

% Latex Beamer Übersicht: 
% https://www2.informatik.hu-berlin.de/~mischulz/beamer.html

\title{WebPiraten}  
\author{A. Flohr, C. Nehls, M. Rasch,  M. Neumann}
\date{22.08.2014} 

\begin{frame}
\maketitle
\end{frame} 

%\begin{frame}
%\frametitle{Inhalt}
%\tableofcontents
%\end{frame}

%\begin{frame}
%\begin{itemize}
%\item Demo
%\item Übersichtgrafik
%\item Use-Case (grob)\newline
%-$>$ GUI -(websocket)-$>$ Preprocossor -(TCP)-$>$ VM -(STDIN)-$>$ Ausführung -(STDIN)-$>$ %VM-Wrapper -(TCP)-$>$ Rails-Simulation -(websocket)-$>$ GUI $<$-
%\item Preprocessor
%\item VM
%\item Simulation
%\item Client
%\item Erlang
%\item Scopes
%\item Sicherheit
%\item Performance
%\end{itemize}
%\end{frame}

\section{Die WebApp}

\section{Use-Case}

\section{Module}
\subsection{Präprozesssor für Ruby}

\begin{frame}
\begin{center}
	\includegraphics[scale=0.35]{overview}
\end{center}
\end{frame}

\begin{frame}
\frametitle{Idee}
%Ruby, rub, rub! :3
\begin{center}
\includegraphics[scale=0.5]{preprocessor/pics/Idee}
\end{center}
%Spiellogik vor den User code kopieren und den Code für das line highlighting und das Verfolgen von Variablen mit dem nötigen Input füttern.
\end{frame}

\begin{frame}
\frametitle{Erkennung von Sonderfällen}
	\includegraphics[scale=0.70]{preprocessor/pics/CodeArten}
\end{frame}

\begin{frame}
\frametitle{Standard Verfahren}
\inputminted[linenos, numbersep=3pt, tabsize=2, frame=lines]{ruby}{preprocessor/sampleStandard.rb}
\end{frame}

\begin{frame}
\frametitle{Standard Verfahren mit Variablenverfolgung}
  \inputminted[linenos, numbersep=5pt, tabsize=4, frame=lines]{ruby}{preprocessor/sampleDebug.rb}
\end{frame}

\begin{frame}
\frametitle{Verarbeitung mehrzeiliger Strings}
  \includegraphics[scale=1]{preprocessor/pics/MultilineVerarbeitung}
\end{frame}

%\begin{frame}
%\frametitle{Verarbeitung von Case Blöcken}
%\begin{center}
%\includegraphics[scale=0.8]{preprocessor/pics/CaseBlockVerarbeitung}
%\end{center}
%\end{frame}

\begin{frame}
%\frametitle{Code: vorher und nachher}
  \begin{columns}
    \column[t]{0.50\textwidth}
\inputminted[linenos, frame=lines, tabsize=2, fontsize=\footnotesize , label=case when einzeilig]{ruby}{preprocessor/CaseA.rb}
    \column[t]{0.50\textwidth}
    \inputminted[linenos, frame=lines, tabsize=2, fontsize=\footnotesize , label=case mehrzeilig]{ruby}{preprocessor/CaseB.rb}
  \end{columns} 
\end{frame}

\begin{frame}
%\frametitle{Code: vorher und nachher}
  \begin{columns}
    \column[t]{.50\textwidth}
\inputminted[ tabsize=2, fontsize=\footnotesize , label=vorher]{ruby}{preprocessor/sampleA.rb}
    \column[t]{.50\textwidth}
    \inputminted[ tabsize=2, fontsize=\footnotesize , label=danach]{ruby}{preprocessor/sampleB.rb}
  \end{columns} 
\end{frame}

\begin{frame}
\frametitle{Beschränkungen für den Nutzer}
  \begin{columns}
    \column[c]{.65\textwidth}
    \begin{itemize}
    \item \textit{return} bei eigenen Funktionen
    \item keine mehrzeiligen Strings der Form             
    \%q[Text des Strings] oder 
    
    \%Q[anderer Text]
    \end{itemize}
    \column[c]{.35\textwidth}
  \includegraphics[scale=0.15]{preprocessor/pics/Verbotsschild}
  \end{columns} 
\end{frame}
\subsection{Codeausführung}

\begin{frame}
\frametitle{API}
\inputminted[linenos, numbersep=2pt, tabsize=4, frame=lines, label=move]{ruby}{vm/move.rb}
\end{frame}

\begin{frame}
\frametitle{API}
\inputminted[linenos, numbersep=2pt, tabsize=4, frame=lines, label=look]{ruby}{vm/look.rb}
\end{frame}

\begin{frame}
\frametitle{VM}
\begin{center}
\includegraphics[width=250pt]{vm/vm.pdf}
\end{center}
\end{frame}
\section{Simulation}

\begin{frame}
Simuliere Krankheit .... DONE
\end{frame}
\subsection{Javascript-Client}

\begin{frame}
\begin{center}
	\includegraphics[scale=0.35]{overview}
\end{center}
\end{frame}

\begin{frame}
\frametitle{Javascript-Client}
\begin{center}
	\includegraphics[scale=0.3]{client/modules.png}
\end{center}
\end{frame}

\begin{frame}
\frametitle{CodeGUI}
\begin{center}
	\includegraphics[scale=0.35]{client/client}
\end{center}
\end{frame}

\begin{frame}
\frametitle{SocketHandler}
	\includegraphics[scale=0.37]{client/socket-queue.PNG}
\end{frame}

\begin{frame}
\frametitle{SocketHandler}
\inputminted[linenos, numbersep=2pt, tabsize=4, frame=lines, label=Beispiel Paket]{json}{client/packet.json}
\end{frame}

\section{Individuelle Arbeiten}

\subsection{Einbindung für Erlang}

\begin{frame}
\begin{center}
\includegraphics[scale=0.5]{erlang/pics/erlang.png}
\end{center}
\end{frame}

\begin{frame}

\end{frame}

\begin{frame}
erlangstuff momentan erhältich
naiver ansatz 

Finde $->$ und kopiere dahinter die line function

%Bild? vllt code
%vorher nachher... neee
\end{frame}

\begin{frame}
%wie bei ruby schiffslogik wird davor kopiert
%und der code entsprechend verarbeitet, \\
laufzeitfehlermeldungen ohne shell hässlich, werden abgefangen und hübsch weitergeleitet
\end{frame}

\begin{frame}
Ziel\\
für besseres highlighting die line-Funktion mit in die vordefinierten Funktionen kopieren
funktionsvariablen verfolgen lassen\\
falsche pfeile erkennen, bzw richtige\\
Diagramm
\end{frame}

\subsection{Scopes}

\begin{frame}
scoobidoobidoo!
\end{frame}
\subsection{Sicherheit}

\begin{frame}
beeedobeeeedobeeeeeeedo
\end{frame}
\subsection{Performance}

\begin{frame}
\begin{center}
\frametitle {Performance Analyse}
\includegraphics[scale=0.4]{performance/clock}
\end{center}
\end{frame}

\begin{frame}
\frametitle{Rechner Phase 1}
\begin{itemize}
	\item Chevalblanc
	\item Virtuelle Maschine (QEMU 1.1.2 mit KVM Support)
	\item 1 CPU Kern - 2.00 GHz (64 bit)
	\item 2 GB RAM
	
\end{itemize}
\end{frame}

\begin{frame}
\frametitle{Websocket-Verbindungen}
	\begin{center}
		\includegraphics[scale=0.3]{performance/connection-response.png}
		% 500 gleichzeitig eingeleitete Verbindungen
		% 500 maximum des Laptops
		% maximale CPU Auslastung von etwa 15%
		% Reiner Verbindungsaufbau
		% Websocket idle bei 500 Verbindungen ist bei < 1% CPU
		% sequentielle Abarbeitung der Verbindungen
		% => Mehr User = linearer Latenzanstieg
		% 1 Thread behandelt websockets
		% => Multithreading (start simulation etc)
		%TODO Berechne durchschnittssteigerung
		%TODO Einfluss der Simulation
	
	\end{center}
\end{frame}

\begin{frame}
\frametitle{StdIn vs StdOut}
\begin{center}
	\includegraphics[scale=0.3]{performance/100-turns.png}
	\includegraphics[scale=0.3]{performance/100-looks.png}
	% 100*look
	% 80ms pro look
	% turn: 0.007824740000000002ms
	% Preprozessor: 2.79% (6.48 ms), Pendelt zwischen 6 und 12ms
\end{center}

\end{frame}

\begin{frame}
	
	
	\begin{center}
	
	\begin{tabular}{c|c}
	Parallele Instanzen & Operationen\\
	\hline% countOps, Beeinflussung
	 1	&562\\
	5	&562\\
	6	&564\\
	7	&567\\
	8	&557\\
	9	&549\\
	10	&545\\
	50&	300\\
	100 &	228\\
	\end{tabular}
	
	\end{center}
	

\end{frame}

\begin{frame}	
	\begin{center}
	
	\begin{tabular}{c|c}
	Parallele Instanzen & CPU Auslastung\\
	\hline
	 1	& 8\%\\
	5	& 28\%\\
	50	& 65\%\\
	100	& 60\%\\
	200	& 46\%\\
	\end{tabular}
	
	\end{center}
\end{frame}

\begin{frame}
\begin{center}
\begin{scriptsize}
\LARGE Ausblick \& Fazit

\end{scriptsize}
\end{center}
\end{frame}

\end{document}